\documentclass[11pt,a4paper]{article}
\usepackage[utf8]{inputenc}
\usepackage[T1]{fontenc}
\usepackage{geometry}
\usepackage{graphicx}
\usepackage{amsmath}
\usepackage{url}
\geometry{margin=2.2cm}
% Reproducible PDF metadata (reduces run-to-run hash drift)
\pdfinfoomitdate=1
\pdftrailerid{}
\pdfsuppressptexinfo=15

\title{ERC8004 Empirics -- Step 06 Top-client address deep profile}
\author{ERC8004-Specialist}
\date{Rebuild timestamp tracked in overnight\_maintenance\_log.md}

\begin{document}
\maketitle

\section*{Executive summary}
Questo step costruisce un profilo completo e mirato del top client individuato in Step 03:
\texttt{0xf653...c807e} (address completo nelle tabelle Step06).
L'analisi copre: composizione semantica dei tag, distribuzione dei valori e score normalizzati,
concentrazione per agent target, dinamica temporale e controlli su range/decimali.

Risultati chiave:
\begin{itemize}
  \item feedback emessi dal top client: \textbf{1215} su 2782 totali (\textbf{43.67\%});
  \item agent distinti toccati: \textbf{1102} su 1470 con almeno un feedback (\textbf{74.97\%});
  \item composizione tag estremamente concentrata in due assi: \textbf{trust} (67.33\%) e \textbf{liveness} (32.67\%);
  \item nessun out-of-range nello score normalizzato: \textbf{0/1215};
  \item nessun uso di decimali su questo address (\texttt{valueDecimals>0}: \textbf{0}).
\end{itemize}

\section{Obiettivo analitico}
L'obiettivo non è solo quantificare il peso volumetrico del top client, ma caratterizzare la sua \emph{firma comportamentale}:
\begin{enumerate}
  \item \textbf{firma semantica} (tag1/tag2 usati),
  \item \textbf{firma numerica} (value, valueDecimals, score normalizzato),
  \item \textbf{firma strutturale} (dispersione dei feedback su agent),
  \item \textbf{firma temporale} (quando si concentra l'attività).
\end{enumerate}

\section{Dati, definizioni e pipeline}
\textbf{Input:}
\begin{itemize}
  \item \path{working/data/raw/2026-03-07_1220_eth_validation_fix/ethereum_1/}
  \item \path{reputation_newfeedback.csv}
\end{itemize}

\textbf{Definizioni operative:}
\begin{itemize}
  \item top client: address con massimo numero di feedback nel dataset; in caso di parità si applica tie-break deterministico (ordinamento lessicografico crescente dell'address) per garantire riproducibilità;
  \item ranking top-client condizionato allo snapshot: cambiando finestra/estrazione input, possono cambiare identity del top client e percentuali condizionate;
  \item score normalizzato:
  \[
  score\_scaled = \frac{value}{10^{valueDecimals}};
  \]
  \item out-of-range: record con \(score\_scaled<0\) oppure \(score\_scaled>100\);
  \item bin temporali per timeline: ampiezza fissa 5000 blocchi.
\end{itemize}

\section{Risultati quantitativi principali}
\subsection*{S1. Sintesi metrica}
\begin{center}
\begin{tabular}{|p{10.4cm}|r|}
\hline
\textbf{Indicatore} & \textbf{Valore} \\
\hline
Feedback top client & 1215 \\
\hline
Share sul totale feedback dataset & 43.67\% \\
\hline
Agent distinti toccati & 1102 \\
\hline
Copertura su agent con almeno un feedback & 74.97\% \\
\hline
Score validi (normalizzati) & 1215 \\
\hline
Score out-of-range [0,100] & 0 (0.00\%) \\
\hline
Record con valueDecimals $>0$ & 0 \\
\hline
Record con decimali significativi post-scaling & 0 \\
\hline
\end{tabular}
\end{center}

\subsection*{S2. Profilo tag}
La firma semantica del top client è molto netta.
Distribuzione \texttt{tag1}:
\begin{itemize}
  \item \texttt{trust}: 818 feedback (67.33\%)
  \item \texttt{liveness}: 397 feedback (32.67\%)
\end{itemize}
Distribuzione \texttt{tag2}:
\begin{itemize}
  \item \texttt{oracle-screening}: 818 (67.33\%)
  \item \texttt{liveness-check}: 397 (32.67\%)
\end{itemize}

Quindi questo address non è ``generalista'': opera quasi esclusivamente su una coppia semantica trust/liveness.
\textbf{Caveat:} tutte le percentuali in questa sezione sono condizionate ai soli 1215 feedback del top client (non alla distribuzione globale del dataset).

\subsection*{S3. Concentrazione su agent target}
Pur avendo volume alto, il top client non è concentrato su pochi agent.
Dalla tabella agent-level:
\begin{itemize}
  \item massimo feedback su singolo agent: 3;
  \item agent con almeno 2 feedback: 105 (su 1102);
  \item quota cumulata top-5 agent: 1.23\%;
  \item quota cumulata top-20 agent: 3.95\%.
\end{itemize}
Interpretazione: strategia \emph{wide coverage}, non \emph{repetition-heavy} su pochi target.
\textbf{Caveat inferenziale:} una dispersione ampia sui target non implica automaticamente assenza di coordinamento; per distinguere screening legittimo da comportamento orchestrato serve integrazione con segnali esterni (es. co-movimento temporale con altri client e metadata off-chain).

\section{Figure (costruzione e lettura)}
\subsection*{F6a -- Distribuzione tag1 del top client}
\textbf{Path:} \texttt{working/results/empirics/figures/fig06a\_top\_client\_tag1\_distribution.pdf}

Bar chart orizzontale dei tag1 ordinati per frequenza (top-15; in pratica 2 etichette non nulle).
\textbf{Caveat di lettura:} la figura misura composizione interna del top client; non va interpretata come quota tag sull'intero universo ERC8004. Con una cardinalità di categorie così bassa, è buona pratica riportare sempre anche i conteggi assoluti (oltre alle percentuali) per evitare sovra-interpretazioni su differenze marginali.

\begin{center}
\includegraphics[width=0.86\textwidth]{../figures/fig06a_top_client_tag1_distribution.pdf}
\end{center}

\subsection*{F6b -- Distribuzione score del top client}
\textbf{Path:} \texttt{working/results/empirics/figures/fig06b\_top\_client\_score\_distribution.pdf}

Istogramma dei soli score in range [0,100], con box informativo su numerosità valida e out-of-range.
Nel caso specifico non ci sono out-of-range.
\textbf{Nota metodologica:} la lettura va fatta congiuntamente a \texttt{valueDecimals}: quando i decimali sono nulli (come qui), la distribuzione riflette direttamente la scala intera originale senza smoothing da normalizzazione frazionaria. Assenza di out-of-range non implica da sola calibrazione ``corretta'': indica solo rispetto del vincolo meccanico di range e va interpretata insieme a dispersione, code e confronto con benchmark esterni.

\begin{center}
\includegraphics[width=0.86\textwidth]{../figures/fig06b_top_client_score_distribution.pdf}
\end{center}

\subsection*{F6c -- Curva di concentrazione per rank agent}
\textbf{Path:} \texttt{working/results/empirics/figures/fig06c\_top\_client\_agent\_concentration.pdf}

Asse x in scala log (rank agent per frequenza), curva principale dei conteggi e curva secondaria di quota cumulata.
Serve per diagnosticare rapidamente se il client è focalizzato su pochi target o distribuito su molti.
\textbf{Caveat di lettura:} con asse rank logaritmico, piccole distanze visive nelle prime posizioni possono corrispondere a differenze assolute rilevanti nei conteggi. Inoltre, quando il massimo feedback per singolo agent è basso (qui: 3), la curva va letta soprattutto come indicatore di ampiezza copertura, non come evidenza di targeting ripetitivo ad alta intensità.

\begin{center}
\includegraphics[width=0.88\textwidth]{../figures/fig06c_top_client_agent_concentration.pdf}
\end{center}

\subsection*{F6d -- Timeline attività (bin 5000 blocchi)}
\textbf{Path:} \texttt{working/results/empirics/figures/fig06d\_top\_client\_timeline\_bins5000.pdf}

Distribuzione temporale dei feedback del top client lungo l'intervallo osservato.
\textbf{Nota metodologica:} i bin da 5000 blocchi privilegiano robustezza/leggibilità rispetto alla granularità intra-bin; picchi isolati vanno interpretati come finestre aggregate e non come burst puntuali a livello transazione. Inoltre, la larghezza in blocchi non è equivalente a finestre temporali uniformi in ore/giorni: la mappatura in tempo reale richiede i timestamp on-chain dei blocchi.

\begin{center}
\includegraphics[width=0.88\textwidth]{../figures/fig06d_top_client_timeline_bins5000.pdf}
\end{center}

\section{Discussione}
Il top client domina il volume complessivo dei feedback, ma lo fa con una dinamica selettiva nei contenuti e diffusa nei target.
Dal lato semantico, l'attività è quasi interamente confinata su trust/liveness;
 dal lato strutturale, i feedback sono dispersi su una platea molto ampia di agent, con bassa ripetizione per coppia client-agent.
Questo pattern è coerente con un comportamento di screening esteso e non con campagne focalizzate su pochi agent.

\textbf{Caveat metodologico (selezione ex-post):} l'address analizzato è definito come massimo osservato nel campione (argmax dei conteggi). Di conseguenza, differenze ``estreme'' rispetto a un client medio possono riflettere anche effetto di selezione. Per inferenze causali/comparative serve confronto con benchmark distribution (es. top-k, quantili, o placebo-rank) e non solo con statistiche puntuali del singolo top client.

Sul piano della qualità numerica, l'assenza di out-of-range e di decimali sul top client rende il suo profilo stabile rispetto al vincolo di score in [0,100].

\section{Output prodotti (Step 06)}
\textbf{Script:}
\begin{itemize}
  \item \texttt{working/analysis/empirics/empirics\_step06\_top\_client\_address.py}
\end{itemize}

\textbf{Tabelle:}
\begin{itemize}
  \item \texttt{working/results/empirics/tables/step06\_top\_client\_profile\_metrics.csv}
  \item \texttt{working/results/empirics/tables/step06\_top\_client\_tag1\_distribution.csv}
  \item \texttt{working/results/empirics/tables/step06\_top\_client\_tag2\_distribution.csv}
  \item \texttt{working/results/empirics/tables/step06\_top\_client\_value\_stats\_by\_tag1.csv}
  \item \texttt{working/results/empirics/tables/step06\_top\_client\_feedback\_by\_agent.csv}
  \item \texttt{working/results/empirics/tables/step06\_top\_client\_time\_bins\_5000.csv}
  \item \texttt{working/results/empirics/tables/step06\_top\_client\_score\_out\_of\_range\_rows.csv}
\end{itemize}

\textbf{Figure (PDF):}
\begin{itemize}
  \item \texttt{working/results/empirics/figures/fig06a\_top\_client\_tag1\_distribution.pdf}
  \item \texttt{working/results/empirics/figures/fig06b\_top\_client\_score\_distribution.pdf}
  \item \texttt{working/results/empirics/figures/fig06c\_top\_client\_agent\_concentration.pdf}
  \item \texttt{working/results/empirics/figures/fig06d\_top\_client\_timeline\_bins5000.pdf}
\end{itemize}

\textbf{Summary breve:}
\begin{itemize}
  \item \texttt{working/results/empirics/summaries/step06\_top\_client\_address\_profile.md}
\end{itemize}

\end{document}
